% figs_part2.tex —— 由 mit/make_figs_part2.py 產生，請勿手動修改（附錄 A.5–A.8，動畫 M05–M08）。
% 重新產生：python notes/mit/make_figs_part2.py

\newcommand{\FigMfiveTrees}{%
\begin{tikzpicture}
\draw[cA,line width=1.4pt,line cap=round] (0,0) -- (-0.475,-0.75);
\draw[cA,line width=1.4pt,line cap=round] (-0.475,-0.75) -- (-0.475,-1.5);
\draw[green!50!black,line width=1.4pt,line cap=round] (0,0) -- (0.475,-0.75);
\fill[black] (0,0) circle (0.085);
\fill[cA] (-0.475,-0.75) circle (0.07);
\fill[cA] (-0.475,-1.5) circle (0.07);
\fill[green!50!black] (0.475,-0.75) circle (0.07);
\node[font=\scriptsize] at (0,0.3) {根};
\draw[green!50!black,line width=1.4pt,line cap=round] (2.6,0) -- (2.125,-0.75);
\draw[cA,line width=1.4pt,line cap=round] (2.6,0) -- (3.075,-0.75);
\draw[cA,line width=1.4pt,line cap=round] (3.075,-0.75) -- (3.075,-1.5);
\fill[black] (2.6,0) circle (0.085);
\fill[green!50!black] (2.125,-0.75) circle (0.07);
\fill[cA] (3.075,-0.75) circle (0.07);
\fill[cA] (3.075,-1.5) circle (0.07);
\node[font=\scriptsize] at (2.6,0.3) {根};
\node[font=\Large] at (1.3,-0.55) {$\neq$};
\node[font=\scriptsize,align=center,text=gray!40!black] at (1.3,-1.95) {不管順序是同一棵樹；\\作為平面樹是不同的兩棵};
\draw[gray!40,line width=0.4pt] (3.95,0.45) -- (3.95,-2.75);
\node[anchor=west,font=\small] at (4.35,0.26) {$n=0$};
\fill[black] (5.85,0.38) circle (0.06);
\node[font=\small] at (11.55,0.26) {$C_0=1$};
\node[anchor=west,font=\small] at (4.35,-0.24) {$n=1$};
\draw[gray!70,line width=0.8pt,line cap=round] (5.85,-0.12) -- (5.85,-0.42);
\fill[black] (5.85,-0.12) circle (0.06);
\fill[cA] (5.85,-0.42) circle (0.045);
\node[font=\small] at (11.55,-0.24) {$C_1=1$};
\node[anchor=west,font=\small] at (4.35,-0.87) {$n=2$};
\draw[gray!70,line width=0.8pt,line cap=round] (5.85,-0.75) -- (5.7,-1.05);
\draw[gray!70,line width=0.8pt,line cap=round] (5.85,-0.75) -- (6,-1.05);
\fill[black] (5.85,-0.75) circle (0.06);
\fill[cA] (5.7,-1.05) circle (0.045);
\fill[cA] (6,-1.05) circle (0.045);
\draw[gray!70,line width=0.8pt,line cap=round] (6.97,-0.75) -- (6.97,-1.05);
\draw[gray!70,line width=0.8pt,line cap=round] (6.97,-1.05) -- (6.97,-1.35);
\fill[black] (6.97,-0.75) circle (0.06);
\fill[cA] (6.97,-1.05) circle (0.045);
\fill[cA] (6.97,-1.35) circle (0.045);
\node[font=\small] at (11.55,-0.87) {$C_2=2$};
\node[anchor=west,font=\small] at (4.35,-1.77) {$n=3$};
\draw[gray!70,line width=0.8pt,line cap=round] (5.85,-1.65) -- (5.55,-1.95);
\draw[gray!70,line width=0.8pt,line cap=round] (5.85,-1.65) -- (5.85,-1.95);
\draw[gray!70,line width=0.8pt,line cap=round] (5.85,-1.65) -- (6.15,-1.95);
\fill[black] (5.85,-1.65) circle (0.06);
\fill[cA] (5.55,-1.95) circle (0.045);
\fill[cA] (5.85,-1.95) circle (0.045);
\fill[cA] (6.15,-1.95) circle (0.045);
\node[font=\tiny,text=gray!60!black] at (5.85,-2.77) {1};
\draw[gray!70,line width=0.8pt,line cap=round] (6.97,-1.65) -- (6.82,-1.95);
\draw[gray!70,line width=0.8pt,line cap=round] (6.97,-1.65) -- (7.12,-1.95);
\draw[gray!70,line width=0.8pt,line cap=round] (7.12,-1.95) -- (7.12,-2.25);
\fill[black] (6.97,-1.65) circle (0.06);
\fill[cA] (6.82,-1.95) circle (0.045);
\fill[cA] (7.12,-1.95) circle (0.045);
\fill[cA] (7.12,-2.25) circle (0.045);
\node[font=\tiny,text=gray!60!black] at (6.97,-2.77) {2};
\draw[gray!70,line width=0.8pt,line cap=round] (8.09,-1.65) -- (7.94,-1.95);
\draw[gray!70,line width=0.8pt,line cap=round] (7.94,-1.95) -- (7.94,-2.25);
\draw[gray!70,line width=0.8pt,line cap=round] (8.09,-1.65) -- (8.24,-1.95);
\fill[black] (8.09,-1.65) circle (0.06);
\fill[cA] (7.94,-1.95) circle (0.045);
\fill[cA] (7.94,-2.25) circle (0.045);
\fill[cA] (8.24,-1.95) circle (0.045);
\node[font=\tiny,text=gray!60!black] at (8.09,-2.77) {3};
\draw[gray!70,line width=0.8pt,line cap=round] (9.21,-1.65) -- (9.21,-1.95);
\draw[gray!70,line width=0.8pt,line cap=round] (9.21,-1.95) -- (9.06,-2.25);
\draw[gray!70,line width=0.8pt,line cap=round] (9.21,-1.95) -- (9.36,-2.25);
\fill[black] (9.21,-1.65) circle (0.06);
\fill[cA] (9.21,-1.95) circle (0.045);
\fill[cA] (9.06,-2.25) circle (0.045);
\fill[cA] (9.36,-2.25) circle (0.045);
\node[font=\tiny,text=gray!60!black] at (9.21,-2.77) {4};
\draw[gray!70,line width=0.8pt,line cap=round] (10.33,-1.65) -- (10.33,-1.95);
\draw[gray!70,line width=0.8pt,line cap=round] (10.33,-1.95) -- (10.33,-2.25);
\draw[gray!70,line width=0.8pt,line cap=round] (10.33,-2.25) -- (10.33,-2.55);
\fill[black] (10.33,-1.65) circle (0.06);
\fill[cA] (10.33,-1.95) circle (0.045);
\fill[cA] (10.33,-2.25) circle (0.045);
\fill[cA] (10.33,-2.55) circle (0.045);
\node[font=\tiny,text=gray!60!black] at (10.33,-2.77) {5};
\node[font=\small] at (11.55,-1.77) {$C_3=5$};
\draw[gray!55,dashed,line width=0.5pt,rounded corners=2pt] (6.55,-1.48) rectangle (8.51,-2.65);
\draw[cTop,line width=0.7pt,rounded corners=2pt] (8.79,-1.48) rectangle (9.63,-2.65);
\end{tikzpicture}}

\newcommand{\FigMsixCut}{%
\begin{tikzpicture}
\draw[gray!65,line width=0.5pt] (0,0) circle (0.9);
\draw[cA,line width=1.1pt,line cap=round] (0,0.9) -- (-0.779,0.45);
\draw[cA,line width=1.1pt,line cap=round] (0.779,0.45) -- (0.779,-0.45);
\draw[cA,line width=1.1pt,line cap=round] (0,-0.9) -- (-0.779,-0.45);
\fill[black] (0,0.9) circle (0.035);
\fill[black] (0.779,0.45) circle (0.035);
\fill[black] (0.779,-0.45) circle (0.035);
\fill[black] (0,-0.9) circle (0.035);
\fill[black] (-0.779,-0.45) circle (0.035);
\fill[black] (-0.779,0.45) circle (0.035);
\node[font=\scriptsize] at (0,1.14) {1};
\node[font=\scriptsize] at (0.987,0.57) {2};
\node[font=\scriptsize] at (0.987,-0.57) {3};
\node[font=\scriptsize] at (0,-1.14) {4};
\node[font=\scriptsize] at (-0.987,-0.57) {5};
\node[font=\scriptsize] at (-0.987,0.57) {6};
\draw[cBad,line width=1.2pt] (-0.412,0.714) -- (-0.55,0.953);
\node[font=\scriptsize,text=cBad,anchor=south east] at (-0.53,0.918) {剪開};
\draw[-{Stealth[length=5pt]},gray!60!black,line width=0.7pt] (1.35,0) -- (2.35,0) node[midway,above,font=\scriptsize] {攤平};
\draw[gray!60,line width=0.4pt] (2.63,-0.45) -- (5.37,-0.45);
\draw[cA,line width=1.1pt] (2.75,-0.45) arc[start angle=180,end angle=0,x radius=1.25,y radius=1.25];
\draw[cA,line width=1.1pt] (3.25,-0.45) arc[start angle=180,end angle=0,x radius=0.25,y radius=0.25];
\draw[cA,line width=1.1pt] (4.25,-0.45) arc[start angle=180,end angle=0,x radius=0.25,y radius=0.25];
\fill (2.75,-0.45) circle (0.03);
\fill (3.25,-0.45) circle (0.03);
\fill (3.75,-0.45) circle (0.03);
\fill (4.25,-0.45) circle (0.03);
\fill (4.75,-0.45) circle (0.03);
\fill (5.25,-0.45) circle (0.03);
\node[font=\scriptsize,inner sep=0pt] at (2.75,-0.65) {1};
\node[font=\scriptsize,inner sep=0pt] at (3.25,-0.65) {2};
\node[font=\scriptsize,inner sep=0pt] at (3.75,-0.65) {3};
\node[font=\scriptsize,inner sep=0pt] at (4.25,-0.65) {4};
\node[font=\scriptsize,inner sep=0pt] at (4.75,-0.65) {5};
\node[font=\scriptsize,inner sep=0pt] at (5.25,-0.65) {6};
\draw[gray!40,line width=0.4pt] (5.85,1.15) -- (5.85,-1.15);
\draw[gray!60,line width=0.4pt] (6.28,0.15) -- (7.72,0.15);
\draw[cA,line width=1.1pt] (6.4,0.15) arc[start angle=180,end angle=0,x radius=0.2,y radius=0.17];
\draw[cA,line width=1.1pt] (7.2,0.15) arc[start angle=180,end angle=0,x radius=0.2,y radius=0.17];
\fill (6.4,0.15) circle (0.03);
\fill (6.8,0.15) circle (0.03);
\fill (7.2,0.15) circle (0.03);
\fill (7.6,0.15) circle (0.03);
\node[font=\scriptsize,inner sep=0pt] at (6.4,-0.02) {i};
\node[font=\scriptsize,inner sep=0pt] at (6.8,-0.02) {j};
\node[font=\scriptsize,inner sep=0pt] at (7.2,-0.02) {k};
\node[font=\scriptsize,inner sep=0pt] at (7.6,-0.02) {l};
\node[font=\small,text=cA] at (7,-0.42) {分開};
\node[font=\small] at (7,-0.85) {$(i,j),(k,l)$};
\draw[gray!60,line width=0.4pt] (8.28,0.15) -- (9.72,0.15);
\draw[cBad,line width=1.1pt] (8.4,0.15) arc[start angle=180,end angle=0,x radius=0.4,y radius=0.34];
\draw[cBad,line width=1.1pt] (8.8,0.15) arc[start angle=180,end angle=0,x radius=0.4,y radius=0.34];
\fill (8.4,0.15) circle (0.03);
\fill (8.8,0.15) circle (0.03);
\fill (9.2,0.15) circle (0.03);
\fill (9.6,0.15) circle (0.03);
\node[font=\scriptsize,inner sep=0pt] at (8.4,-0.02) {i};
\node[font=\scriptsize,inner sep=0pt] at (8.8,-0.02) {j};
\node[font=\scriptsize,inner sep=0pt] at (9.2,-0.02) {k};
\node[font=\scriptsize,inner sep=0pt] at (9.6,-0.02) {l};
\node[font=\small,text=cBad] at (9,-0.42) {交叉};
\node[font=\small] at (9,-0.85) {$(i,k),(j,l)$};
\draw[gray!60,line width=0.4pt] (10.28,0.15) -- (11.72,0.15);
\draw[violet!75!black,line width=1.1pt] (10.4,0.15) arc[start angle=180,end angle=0,x radius=0.6,y radius=0.51];
\draw[violet!75!black,line width=1.1pt] (10.8,0.15) arc[start angle=180,end angle=0,x radius=0.2,y radius=0.17];
\fill (10.4,0.15) circle (0.03);
\fill (10.8,0.15) circle (0.03);
\fill (11.2,0.15) circle (0.03);
\fill (11.6,0.15) circle (0.03);
\node[font=\scriptsize,inner sep=0pt] at (10.4,-0.02) {i};
\node[font=\scriptsize,inner sep=0pt] at (10.8,-0.02) {j};
\node[font=\scriptsize,inner sep=0pt] at (11.2,-0.02) {k};
\node[font=\scriptsize,inner sep=0pt] at (11.6,-0.02) {l};
\node[font=\small,text=violet!75!black] at (11,-0.42) {巢狀};
\node[font=\small] at (11,-0.85) {$(i,l),(j,k)$};
\node[font=\small] at (9.0,1.0) {$i<j<k<l$};
\end{tikzpicture}}

\newcommand{\FigMsixCount}{%
\begin{tikzpicture}
\draw[gray!60,line width=0.4pt] (-0.1,0) -- (1.2,0);
\draw[cA,line width=0.9pt] (0,0) arc[start angle=180,end angle=0,x radius=0.11,y radius=0.093];
\draw[cA,line width=0.9pt] (0.44,0) arc[start angle=180,end angle=0,x radius=0.11,y radius=0.094];
\draw[cA,line width=0.9pt] (0.88,0) arc[start angle=180,end angle=0,x radius=0.11,y radius=0.094];
\fill (0,0) circle (0.022);
\fill (0.22,0) circle (0.022);
\fill (0.44,0) circle (0.022);
\fill (0.66,0) circle (0.022);
\fill (0.88,0) circle (0.022);
\fill (1.1,0) circle (0.022);
\draw[cA,line width=0.75pt,rounded corners=2pt] (-0.16,-0.08) rectangle (1.26,0.56);
\draw[green!50!black,dashed,line width=0.75pt,rounded corners=2pt] (-0.23,-0.14) rectangle (1.33,0.62);
\draw[gray!60,line width=0.4pt] (-0.1,-0.86) -- (1.2,-0.86);
\draw[cA,line width=0.9pt] (0,-0.86) arc[start angle=180,end angle=0,x radius=0.11,y radius=0.093];
\draw[cA,line width=0.9pt] (0.44,-0.86) arc[start angle=180,end angle=0,x radius=0.22,y radius=0.187];
\draw[cA,line width=0.9pt] (0.66,-0.86) arc[start angle=180,end angle=0,x radius=0.22,y radius=0.187];
\fill (0,-0.86) circle (0.022);
\fill (0.22,-0.86) circle (0.022);
\fill (0.44,-0.86) circle (0.022);
\fill (0.66,-0.86) circle (0.022);
\fill (0.88,-0.86) circle (0.022);
\fill (1.1,-0.86) circle (0.022);
\draw[green!50!black,dashed,line width=0.75pt,rounded corners=2pt] (-0.23,-1) rectangle (1.33,-0.24);
\draw[gray!60,line width=0.4pt] (-0.1,-1.72) -- (1.2,-1.72);
\draw[cA,line width=0.9pt] (0,-1.72) arc[start angle=180,end angle=0,x radius=0.11,y radius=0.093];
\draw[cA,line width=0.9pt] (0.44,-1.72) arc[start angle=180,end angle=0,x radius=0.33,y radius=0.281];
\draw[cA,line width=0.9pt] (0.66,-1.72) arc[start angle=180,end angle=0,x radius=0.11,y radius=0.093];
\fill (0,-1.72) circle (0.022);
\fill (0.22,-1.72) circle (0.022);
\fill (0.44,-1.72) circle (0.022);
\fill (0.66,-1.72) circle (0.022);
\fill (0.88,-1.72) circle (0.022);
\fill (1.1,-1.72) circle (0.022);
\draw[cA,line width=0.75pt,rounded corners=2pt] (-0.16,-1.8) rectangle (1.26,-1.16);
\draw[gray!60,line width=0.4pt] (1.68,0) -- (2.98,0);
\draw[cA,line width=0.9pt] (1.78,0) arc[start angle=180,end angle=0,x radius=0.22,y radius=0.187];
\draw[cA,line width=0.9pt] (2,0) arc[start angle=180,end angle=0,x radius=0.22,y radius=0.187];
\draw[cA,line width=0.9pt] (2.66,0) arc[start angle=180,end angle=0,x radius=0.11,y radius=0.093];
\fill (1.78,0) circle (0.022);
\fill (2,0) circle (0.022);
\fill (2.22,0) circle (0.022);
\fill (2.44,0) circle (0.022);
\fill (2.66,0) circle (0.022);
\fill (2.88,0) circle (0.022);
\draw[green!50!black,dashed,line width=0.75pt,rounded corners=2pt] (1.55,-0.14) rectangle (3.11,0.62);
\draw[gray!60,line width=0.4pt] (1.68,-0.86) -- (2.98,-0.86);
\draw[cA,line width=0.9pt] (1.78,-0.86) arc[start angle=180,end angle=0,x radius=0.22,y radius=0.187];
\draw[cA,line width=0.9pt] (2,-0.86) arc[start angle=180,end angle=0,x radius=0.33,y radius=0.281];
\draw[cA,line width=0.9pt] (2.44,-0.86) arc[start angle=180,end angle=0,x radius=0.22,y radius=0.187];
\fill (1.78,-0.86) circle (0.022);
\fill (2,-0.86) circle (0.022);
\fill (2.22,-0.86) circle (0.022);
\fill (2.44,-0.86) circle (0.022);
\fill (2.66,-0.86) circle (0.022);
\fill (2.88,-0.86) circle (0.022);
\draw[green!50!black,dashed,line width=0.75pt,rounded corners=2pt] (1.55,-1) rectangle (3.11,-0.24);
\draw[gray!60,line width=0.4pt] (1.68,-1.72) -- (2.98,-1.72);
\draw[cA,line width=0.9pt] (1.78,-1.72) arc[start angle=180,end angle=0,x radius=0.22,y radius=0.187];
\draw[cA,line width=0.9pt] (2,-1.72) arc[start angle=180,end angle=0,x radius=0.44,y radius=0.374];
\draw[cA,line width=0.9pt] (2.44,-1.72) arc[start angle=180,end angle=0,x radius=0.11,y radius=0.094];
\fill (1.78,-1.72) circle (0.022);
\fill (2,-1.72) circle (0.022);
\fill (2.22,-1.72) circle (0.022);
\fill (2.44,-1.72) circle (0.022);
\fill (2.66,-1.72) circle (0.022);
\fill (2.88,-1.72) circle (0.022);
\draw[gray!60,line width=0.4pt] (3.46,0) -- (4.76,0);
\draw[cA,line width=0.9pt] (3.56,0) arc[start angle=180,end angle=0,x radius=0.33,y radius=0.28];
\draw[cA,line width=0.9pt] (3.78,0) arc[start angle=180,end angle=0,x radius=0.11,y radius=0.093];
\draw[cA,line width=0.9pt] (4.44,0) arc[start angle=180,end angle=0,x radius=0.11,y radius=0.093];
\fill (3.56,0) circle (0.022);
\fill (3.78,0) circle (0.022);
\fill (4,0) circle (0.022);
\fill (4.22,0) circle (0.022);
\fill (4.44,0) circle (0.022);
\fill (4.66,0) circle (0.022);
\draw[cA,line width=0.75pt,rounded corners=2pt] (3.4,-0.08) rectangle (4.82,0.56);
\draw[gray!60,line width=0.4pt] (3.46,-0.86) -- (4.76,-0.86);
\draw[cA,line width=0.9pt] (3.56,-0.86) arc[start angle=180,end angle=0,x radius=0.33,y radius=0.28];
\draw[cA,line width=0.9pt] (3.78,-0.86) arc[start angle=180,end angle=0,x radius=0.33,y radius=0.281];
\draw[cA,line width=0.9pt] (4,-0.86) arc[start angle=180,end angle=0,x radius=0.33,y radius=0.281];
\fill (3.56,-0.86) circle (0.022);
\fill (3.78,-0.86) circle (0.022);
\fill (4,-0.86) circle (0.022);
\fill (4.22,-0.86) circle (0.022);
\fill (4.44,-0.86) circle (0.022);
\fill (4.66,-0.86) circle (0.022);
\draw[green!50!black,dashed,line width=0.75pt,rounded corners=2pt] (3.33,-1) rectangle (4.89,-0.24);
\draw[gray!60,line width=0.4pt] (3.46,-1.72) -- (4.76,-1.72);
\draw[cA,line width=0.9pt] (3.56,-1.72) arc[start angle=180,end angle=0,x radius=0.33,y radius=0.28];
\draw[cA,line width=0.9pt] (3.78,-1.72) arc[start angle=180,end angle=0,x radius=0.44,y radius=0.374];
\draw[cA,line width=0.9pt] (4,-1.72) arc[start angle=180,end angle=0,x radius=0.22,y radius=0.187];
\fill (3.56,-1.72) circle (0.022);
\fill (3.78,-1.72) circle (0.022);
\fill (4,-1.72) circle (0.022);
\fill (4.22,-1.72) circle (0.022);
\fill (4.44,-1.72) circle (0.022);
\fill (4.66,-1.72) circle (0.022);
\draw[gray!60,line width=0.4pt] (5.24,0) -- (6.54,0);
\draw[cA,line width=0.9pt] (5.34,0) arc[start angle=180,end angle=0,x radius=0.44,y radius=0.374];
\draw[cA,line width=0.9pt] (5.56,0) arc[start angle=180,end angle=0,x radius=0.11,y radius=0.094];
\draw[cA,line width=0.9pt] (6,0) arc[start angle=180,end angle=0,x radius=0.22,y radius=0.187];
\fill (5.34,0) circle (0.022);
\fill (5.56,0) circle (0.022);
\fill (5.78,0) circle (0.022);
\fill (6,0) circle (0.022);
\fill (6.22,0) circle (0.022);
\fill (6.44,0) circle (0.022);
\draw[gray!60,line width=0.4pt] (5.24,-0.86) -- (6.54,-0.86);
\draw[cA,line width=0.9pt] (5.34,-0.86) arc[start angle=180,end angle=0,x radius=0.44,y radius=0.374];
\draw[cA,line width=0.9pt] (5.56,-0.86) arc[start angle=180,end angle=0,x radius=0.22,y radius=0.187];
\draw[cA,line width=0.9pt] (5.78,-0.86) arc[start angle=180,end angle=0,x radius=0.33,y radius=0.28];
\fill (5.34,-0.86) circle (0.022);
\fill (5.56,-0.86) circle (0.022);
\fill (5.78,-0.86) circle (0.022);
\fill (6,-0.86) circle (0.022);
\fill (6.22,-0.86) circle (0.022);
\fill (6.44,-0.86) circle (0.022);
\draw[gray!60,line width=0.4pt] (5.24,-1.72) -- (6.54,-1.72);
\draw[cA,line width=0.9pt] (5.34,-1.72) arc[start angle=180,end angle=0,x radius=0.44,y radius=0.374];
\draw[cA,line width=0.9pt] (5.56,-1.72) arc[start angle=180,end angle=0,x radius=0.44,y radius=0.374];
\draw[cA,line width=0.9pt] (5.78,-1.72) arc[start angle=180,end angle=0,x radius=0.11,y radius=0.093];
\fill (5.34,-1.72) circle (0.022);
\fill (5.56,-1.72) circle (0.022);
\fill (5.78,-1.72) circle (0.022);
\fill (6,-1.72) circle (0.022);
\fill (6.22,-1.72) circle (0.022);
\fill (6.44,-1.72) circle (0.022);
\draw[gray!60,line width=0.4pt] (7.02,0) -- (8.32,0);
\draw[cA,line width=0.9pt] (7.12,0) arc[start angle=180,end angle=0,x radius=0.55,y radius=0.468];
\draw[cA,line width=0.9pt] (7.34,0) arc[start angle=180,end angle=0,x radius=0.11,y radius=0.094];
\draw[cA,line width=0.9pt] (7.78,0) arc[start angle=180,end angle=0,x radius=0.11,y radius=0.093];
\fill (7.12,0) circle (0.022);
\fill (7.34,0) circle (0.022);
\fill (7.56,0) circle (0.022);
\fill (7.78,0) circle (0.022);
\fill (8,0) circle (0.022);
\fill (8.22,0) circle (0.022);
\draw[cA,line width=0.75pt,rounded corners=2pt] (6.96,-0.08) rectangle (8.38,0.56);
\draw[gray!60,line width=0.4pt] (7.02,-0.86) -- (8.32,-0.86);
\draw[cA,line width=0.9pt] (7.12,-0.86) arc[start angle=180,end angle=0,x radius=0.55,y radius=0.468];
\draw[cA,line width=0.9pt] (7.34,-0.86) arc[start angle=180,end angle=0,x radius=0.22,y radius=0.187];
\draw[cA,line width=0.9pt] (7.56,-0.86) arc[start angle=180,end angle=0,x radius=0.22,y radius=0.187];
\fill (7.12,-0.86) circle (0.022);
\fill (7.34,-0.86) circle (0.022);
\fill (7.56,-0.86) circle (0.022);
\fill (7.78,-0.86) circle (0.022);
\fill (8,-0.86) circle (0.022);
\fill (8.22,-0.86) circle (0.022);
\draw[gray!60,line width=0.4pt] (7.02,-1.72) -- (8.32,-1.72);
\draw[cA,line width=0.9pt] (7.12,-1.72) arc[start angle=180,end angle=0,x radius=0.55,y radius=0.468];
\draw[cA,line width=0.9pt] (7.34,-1.72) arc[start angle=180,end angle=0,x radius=0.33,y radius=0.281];
\draw[cA,line width=0.9pt] (7.56,-1.72) arc[start angle=180,end angle=0,x radius=0.11,y radius=0.093];
\fill (7.12,-1.72) circle (0.022);
\fill (7.34,-1.72) circle (0.022);
\fill (7.56,-1.72) circle (0.022);
\fill (7.78,-1.72) circle (0.022);
\fill (8,-1.72) circle (0.022);
\fill (8.22,-1.72) circle (0.022);
\draw[cA,line width=0.75pt,rounded corners=2pt] (6.96,-1.8) rectangle (8.38,-1.16);
\node[font=\scriptsize,text=gray!50!black] at (0.55,0.93) {點 1 配點 2};
\node[font=\scriptsize,text=gray!50!black] at (2.33,0.93) {點 1 配點 3};
\node[font=\scriptsize,text=gray!50!black] at (4.11,0.93) {點 1 配點 4};
\node[font=\scriptsize,text=gray!50!black] at (5.89,0.93) {點 1 配點 5};
\node[font=\scriptsize,text=gray!50!black] at (7.67,0.93) {點 1 配點 6};
\draw[gray!40,line width=0.4pt] (8.9,1.05) -- (8.9,-2.05);
\node[font=\scriptsize,text=cA] at (9.85,0.93) {不交叉};
\node[font=\scriptsize] at (11.3,0.93) {Dyck 字};
\node[font=\scriptsize,text=green!40!black] at (12.8,0.93) {不巢狀};
\draw[gray!60,line width=0.4pt] (9.27,0.38) -- (10.43,0.38);
\draw[cA,line width=0.9pt] (9.35,0.38) arc[start angle=180,end angle=0,x radius=0.1,y radius=0.075];
\draw[cA,line width=0.9pt] (9.75,0.38) arc[start angle=180,end angle=0,x radius=0.1,y radius=0.075];
\draw[cA,line width=0.9pt] (10.15,0.38) arc[start angle=180,end angle=0,x radius=0.1,y radius=0.075];
\fill (9.35,0.38) circle (0.02);
\fill (9.55,0.38) circle (0.02);
\fill (9.75,0.38) circle (0.02);
\fill (9.95,0.38) circle (0.02);
\fill (10.15,0.38) circle (0.02);
\fill (10.35,0.38) circle (0.02);
\node[font=\scriptsize] at (11.3,0.51) {$\mathrm{UDUDUD}$};
\draw[gray!60,line width=0.4pt] (12.22,0.38) -- (13.38,0.38);
\draw[green!50!black,line width=0.9pt] (12.3,0.38) arc[start angle=180,end angle=0,x radius=0.1,y radius=0.075];
\draw[green!50!black,line width=0.9pt] (12.7,0.38) arc[start angle=180,end angle=0,x radius=0.1,y radius=0.075];
\draw[green!50!black,line width=0.9pt] (13.1,0.38) arc[start angle=180,end angle=0,x radius=0.1,y radius=0.075];
\fill (12.3,0.38) circle (0.02);
\fill (12.5,0.38) circle (0.02);
\fill (12.7,0.38) circle (0.02);
\fill (12.9,0.38) circle (0.02);
\fill (13.1,0.38) circle (0.02);
\fill (13.3,0.38) circle (0.02);
\draw[gray!60,line width=0.4pt] (9.27,-0.22) -- (10.43,-0.22);
\draw[cA,line width=0.9pt] (9.35,-0.22) arc[start angle=180,end angle=0,x radius=0.1,y radius=0.075];
\draw[cA,line width=0.9pt] (9.75,-0.22) arc[start angle=180,end angle=0,x radius=0.3,y radius=0.225];
\draw[cA,line width=0.9pt] (9.95,-0.22) arc[start angle=180,end angle=0,x radius=0.1,y radius=0.075];
\fill (9.35,-0.22) circle (0.02);
\fill (9.55,-0.22) circle (0.02);
\fill (9.75,-0.22) circle (0.02);
\fill (9.95,-0.22) circle (0.02);
\fill (10.15,-0.22) circle (0.02);
\fill (10.35,-0.22) circle (0.02);
\node[font=\scriptsize] at (11.3,-0.09) {$\mathrm{UDUUDD}$};
\draw[gray!60,line width=0.4pt] (12.22,-0.22) -- (13.38,-0.22);
\draw[green!50!black,line width=0.9pt] (12.3,-0.22) arc[start angle=180,end angle=0,x radius=0.1,y radius=0.075];
\draw[green!50!black,line width=0.9pt] (12.7,-0.22) arc[start angle=180,end angle=0,x radius=0.2,y radius=0.15];
\draw[green!50!black,line width=0.9pt] (12.9,-0.22) arc[start angle=180,end angle=0,x radius=0.2,y radius=0.15];
\fill (12.3,-0.22) circle (0.02);
\fill (12.5,-0.22) circle (0.02);
\fill (12.7,-0.22) circle (0.02);
\fill (12.9,-0.22) circle (0.02);
\fill (13.1,-0.22) circle (0.02);
\fill (13.3,-0.22) circle (0.02);
\draw[gray!60,line width=0.4pt] (9.27,-0.82) -- (10.43,-0.82);
\draw[cA,line width=0.9pt] (9.35,-0.82) arc[start angle=180,end angle=0,x radius=0.3,y radius=0.225];
\draw[cA,line width=0.9pt] (9.55,-0.82) arc[start angle=180,end angle=0,x radius=0.1,y radius=0.075];
\draw[cA,line width=0.9pt] (10.15,-0.82) arc[start angle=180,end angle=0,x radius=0.1,y radius=0.075];
\fill (9.35,-0.82) circle (0.02);
\fill (9.55,-0.82) circle (0.02);
\fill (9.75,-0.82) circle (0.02);
\fill (9.95,-0.82) circle (0.02);
\fill (10.15,-0.82) circle (0.02);
\fill (10.35,-0.82) circle (0.02);
\node[font=\scriptsize] at (11.3,-0.69) {$\mathrm{UUDDUD}$};
\draw[gray!60,line width=0.4pt] (12.22,-0.82) -- (13.38,-0.82);
\draw[green!50!black,line width=0.9pt] (12.3,-0.82) arc[start angle=180,end angle=0,x radius=0.2,y radius=0.15];
\draw[green!50!black,line width=0.9pt] (12.5,-0.82) arc[start angle=180,end angle=0,x radius=0.2,y radius=0.15];
\draw[green!50!black,line width=0.9pt] (13.1,-0.82) arc[start angle=180,end angle=0,x radius=0.1,y radius=0.075];
\fill (12.3,-0.82) circle (0.02);
\fill (12.5,-0.82) circle (0.02);
\fill (12.7,-0.82) circle (0.02);
\fill (12.9,-0.82) circle (0.02);
\fill (13.1,-0.82) circle (0.02);
\fill (13.3,-0.82) circle (0.02);
\draw[gray!60,line width=0.4pt] (9.27,-1.42) -- (10.43,-1.42);
\draw[cA,line width=0.9pt] (9.35,-1.42) arc[start angle=180,end angle=0,x radius=0.5,y radius=0.375];
\draw[cA,line width=0.9pt] (9.55,-1.42) arc[start angle=180,end angle=0,x radius=0.1,y radius=0.075];
\draw[cA,line width=0.9pt] (9.95,-1.42) arc[start angle=180,end angle=0,x radius=0.1,y radius=0.075];
\fill (9.35,-1.42) circle (0.02);
\fill (9.55,-1.42) circle (0.02);
\fill (9.75,-1.42) circle (0.02);
\fill (9.95,-1.42) circle (0.02);
\fill (10.15,-1.42) circle (0.02);
\fill (10.35,-1.42) circle (0.02);
\node[font=\scriptsize] at (11.3,-1.29) {$\mathrm{UUDUDD}$};
\draw[gray!60,line width=0.4pt] (12.22,-1.42) -- (13.38,-1.42);
\draw[green!50!black,line width=0.9pt] (12.3,-1.42) arc[start angle=180,end angle=0,x radius=0.2,y radius=0.15];
\draw[green!50!black,line width=0.9pt] (12.5,-1.42) arc[start angle=180,end angle=0,x radius=0.3,y radius=0.225];
\draw[green!50!black,line width=0.9pt] (12.9,-1.42) arc[start angle=180,end angle=0,x radius=0.2,y radius=0.15];
\fill (12.3,-1.42) circle (0.02);
\fill (12.5,-1.42) circle (0.02);
\fill (12.7,-1.42) circle (0.02);
\fill (12.9,-1.42) circle (0.02);
\fill (13.1,-1.42) circle (0.02);
\fill (13.3,-1.42) circle (0.02);
\draw[gray!60,line width=0.4pt] (9.27,-2.02) -- (10.43,-2.02);
\draw[cA,line width=0.9pt] (9.35,-2.02) arc[start angle=180,end angle=0,x radius=0.5,y radius=0.375];
\draw[cA,line width=0.9pt] (9.55,-2.02) arc[start angle=180,end angle=0,x radius=0.3,y radius=0.225];
\draw[cA,line width=0.9pt] (9.75,-2.02) arc[start angle=180,end angle=0,x radius=0.1,y radius=0.075];
\fill (9.35,-2.02) circle (0.02);
\fill (9.55,-2.02) circle (0.02);
\fill (9.75,-2.02) circle (0.02);
\fill (9.95,-2.02) circle (0.02);
\fill (10.15,-2.02) circle (0.02);
\fill (10.35,-2.02) circle (0.02);
\node[font=\scriptsize] at (11.3,-1.89) {$\mathrm{UUUDDD}$};
\draw[gray!60,line width=0.4pt] (12.22,-2.02) -- (13.38,-2.02);
\draw[green!50!black,line width=0.9pt] (12.3,-2.02) arc[start angle=180,end angle=0,x radius=0.3,y radius=0.225];
\draw[green!50!black,line width=0.9pt] (12.5,-2.02) arc[start angle=180,end angle=0,x radius=0.3,y radius=0.225];
\draw[green!50!black,line width=0.9pt] (12.7,-2.02) arc[start angle=180,end angle=0,x radius=0.3,y radius=0.225];
\fill (12.3,-2.02) circle (0.02);
\fill (12.5,-2.02) circle (0.02);
\fill (12.7,-2.02) circle (0.02);
\fill (12.9,-2.02) circle (0.02);
\fill (13.1,-2.02) circle (0.02);
\fill (13.3,-2.02) circle (0.02);
\draw[cTop,line width=0.6pt,rounded corners=2pt] (9.2,-1.51) rectangle (13.45,-0.95);
\end{tikzpicture}}

\newcommand{\FigMsevenDual}{%
\begin{tikzpicture}
\fill[cA!10] (-0.676,-0.93) -- (-1.094,0.355) -- (1.094,0.355) -- cycle;
\fill[cA!10] (-1.094,0.355) -- (0,1.15) -- (1.094,0.355) -- cycle;
\fill[cTop!30] (-0.676,-0.93) -- (1.094,0.355) -- (0.676,-0.93) -- cycle;
\draw[gray!70,line width=0.6pt,line join=round] (-0.676,-0.93) -- (-1.094,0.355) -- (0,1.15) -- (1.094,0.355) -- (0.676,-0.93) -- cycle;
\draw[cA,line width=1.1pt,line cap=round] (-0.676,-0.93) -- (1.094,0.355);
\draw[cA,line width=1.1pt,line cap=round] (-1.094,0.355) -- (1.094,0.355);
\draw[cTop,line width=2pt,line cap=round] (-0.676,-0.93) -- (0.676,-0.93);
\node[font=\scriptsize,text=cTop] at (0,-1.13) {底邊};
\fill (-0.676,-0.93) circle (0.033);
\node[font=\scriptsize] at (-0.805,-1.108) {1};
\fill (-1.094,0.355) circle (0.033);
\node[font=\scriptsize] at (-1.303,0.423) {2};
\fill (0,1.15) circle (0.033);
\node[font=\scriptsize] at (0,1.37) {3};
\fill (1.094,0.355) circle (0.033);
\node[font=\scriptsize] at (1.303,0.423) {4};
\fill (0.676,-0.93) circle (0.033);
\node[font=\scriptsize] at (0.805,-1.108) {5};
\draw[black!50,line width=0.8pt,line cap=round] (-0.225,-0.073) -- (-1.227,-0.399);
\draw[black!50,line width=0.8pt,line cap=round] (-0.225,-0.073) -- (0,0.62);
\draw[black!50,line width=0.8pt,line cap=round] (0,0.62) -- (-0.758,1.044);
\draw[black!50,line width=0.8pt,line cap=round] (0,0.62) -- (0.758,1.044);
\draw[black!50,line width=0.8pt,line cap=round] (0.365,-0.502) -- (-0.225,-0.073);
\draw[black!50,line width=0.8pt,line cap=round] (0.365,-0.502) -- (1.227,-0.399);
\fill[gray!60] (-1.269,-0.441) rectangle (-1.185,-0.357);
\fill[gray!60] (-0.8,1.002) rectangle (-0.716,1.086);
\fill[gray!60] (0.716,1.002) rectangle (0.8,1.086);
\fill[gray!60] (1.185,-0.441) rectangle (1.269,-0.357);
\fill[cA] (-0.225,-0.073) circle (0.055);
\node[font=\scriptsize] at (-0.385,0.077) {B};
\fill[cA] (0,0.62) circle (0.055);
\node[font=\scriptsize] at (0.17,0.52) {C};
\fill[cTop] (0.365,-0.502) circle (0.055);
\node[font=\scriptsize] at (0.385,-0.702) {A};
\draw[-{Stealth[length=5pt]},gray!60!black,line width=0.7pt] (1.85,0.05) -- (3.05,0.05) node[midway,above,font=\scriptsize] {拿掉多邊形};
\draw[cLeft,line width=1.0pt,line cap=round] (3.85,0.141) -- (3.35,-0.359);
\draw[cLeft,line width=1.0pt,line cap=round] (3.85,0.141) -- (4.1,-0.359);
\draw[cLeft,line width=1.0pt,line cap=round] (4.1,-0.359) -- (3.85,-0.859);
\draw[cLeft,line width=1.0pt,line cap=round] (4.1,-0.359) -- (4.35,-0.859);
\draw[cTop,line width=1.0pt,line cap=round] (4.1,0.641) -- (3.85,0.141);
\draw[cTop,line width=1.0pt,line cap=round] (4.1,0.641) -- (4.85,0.141);
\fill[cLeft] (3.85,0.141) circle (0.055);
\fill[cLeft] (4.1,-0.359) circle (0.055);
\fill[cTop] (4.1,0.641) circle (0.055);
\fill[cLeft] (3.308,-0.401) rectangle (3.392,-0.317);
\fill[cLeft] (3.808,-0.901) rectangle (3.892,-0.817);
\fill[cLeft] (4.308,-0.901) rectangle (4.392,-0.817);
\fill[cRight] (4.808,0.099) rectangle (4.892,0.183);
\node[font=\scriptsize,inner sep=0pt] at (3.67,0.201) {B};
\node[font=\scriptsize,inner sep=0pt] at (4.28,-0.299) {C};
\node[font=\scriptsize,inner sep=0pt] at (4.1,0.831) {A};
\draw[gray!40,line width=0.4pt] (5.5,1.3) -- (5.5,-1.3);
\draw[cLeft,line width=1.2pt,line cap=round] (6.55,0.55) -- (6.05,-0.15);
\draw[cRight,line width=1.2pt,line cap=round] (6.05,-0.15) -- (6.55,-0.85);
\fill[cA] (6.55,0.55) circle (0.06);
\node[font=\scriptsize] at (6.73,0.6) {A};
\fill[cA] (6.05,-0.15) circle (0.06);
\node[font=\scriptsize] at (5.87,-0.15) {B};
\fill[cA] (6.55,-0.85) circle (0.06);
\node[font=\scriptsize] at (6.73,-0.85) {C};
\node[font=\tiny,text=cLeft] at (6.15,0.33) {左};
\node[font=\tiny,text=cRight] at (6.17,-0.6) {右};
\draw[-{Stealth[length=5pt]},gray!60!black,line width=0.7pt] (7.2,-0.15) -- (9.75,-0.15) node[midway,above,font=\scriptsize,align=center] {\textcolor{cLeft}{左小孩} $\to$ 第一個小孩\\\textcolor{cRight}{右小孩} $\to$ 下一個兄弟} node[midway,below,font=\scriptsize] {加一個新的根};
\draw[gray!60,line width=1.2pt,line cap=round] (10.6,0.75) -- (10.6,0.13);
\draw[cLeft,line width=1.2pt,line cap=round] (10.6,0.13) -- (10.15,-0.49);
\draw[cRight,line width=1.2pt,line cap=round] (10.6,0.13) -- (11.05,-0.49);
\fill[white] (10.6,0.75) circle (0.07);
\draw[black,line width=0.6pt] (10.6,0.75) circle (0.07);
\fill[cA] (10.6,0.13) circle (0.06);
\fill[cA] (10.15,-0.49) circle (0.06);
\fill[cA] (11.05,-0.49) circle (0.06);
\node[font=\scriptsize,inner sep=0pt] at (10.6,0.97) {新的根};
\node[font=\scriptsize,inner sep=0pt] at (10.78,0.17) {A};
\node[font=\scriptsize,inner sep=0pt] at (9.97,-0.49) {B};
\node[font=\scriptsize,inner sep=0pt] at (11.23,-0.49) {C};
\end{tikzpicture}}

\newcommand{\FigMeightWalk}{%
\begin{tikzpicture}
\draw[gray!70,line width=1.1pt,line cap=round] (0,0.35) -- (0,-0.6);
\draw[gray!70,line width=1.1pt,line cap=round] (0,-0.6) -- (-0.8,-1.55);
\draw[gray!70,line width=1.1pt,line cap=round] (0,-0.6) -- (0.8,-1.55);
\fill (0,0.35) circle (0.07);
\fill[cA] (0,-0.6) circle (0.06);
\fill[cA] (-0.8,-1.55) circle (0.06);
\fill[cA] (0.8,-1.55) circle (0.06);
\node[font=\scriptsize] at (0,0.59) {根};
\node[font=\scriptsize] at (0.3,-0.58) {A};
\node[font=\scriptsize] at (-0.8,-1.79) {B};
\node[font=\scriptsize] at (0.8,-1.79) {C};
\node[font=\scriptsize,text=cTop] at (-0.24,-0.125) {1$\downarrow$};
\node[font=\scriptsize,text=cTop] at (0.24,-0.125) {6$\uparrow$};
\node[font=\scriptsize,text=cTop] at (-0.584,-0.92) {2$\downarrow$};
\node[font=\scriptsize,text=cTop] at (-0.216,-1.23) {3$\uparrow$};
\node[font=\scriptsize,text=cTop] at (0.216,-1.23) {4$\downarrow$};
\node[font=\scriptsize,text=cTop] at (0.584,-0.92) {5$\uparrow$};
\draw[gray!55,line width=0.4pt] (1.75,-0.3) -- (4.75,-0.3);
\draw[cA,line width=1.3pt,line cap=round] (1.75,-0.3) -- (2.25,0.2);
\draw[cA,line width=1.3pt,line cap=round] (2.25,0.2) -- (2.75,0.7);
\draw[cA,line width=1.3pt,line cap=round] (2.75,0.7) -- (3.25,0.2);
\draw[cA,line width=1.3pt,line cap=round] (3.25,0.2) -- (3.75,0.7);
\draw[cA,line width=1.3pt,line cap=round] (3.75,0.7) -- (4.25,0.2);
\draw[cA,line width=1.3pt,line cap=round] (4.25,0.2) -- (4.75,-0.3);
\node[font=\scriptsize] at (2,-0.52) {U};
\node[font=\scriptsize] at (2.5,-0.52) {U};
\node[font=\scriptsize] at (3,-0.52) {D};
\node[font=\scriptsize] at (3.5,-0.52) {U};
\node[font=\scriptsize] at (4,-0.52) {D};
\node[font=\scriptsize] at (4.5,-0.52) {D};
\node[font=\scriptsize,anchor=west,text=gray!50!black] at (4.83,-0.3) {地面};
\draw[gray!60,line width=0.4pt] (1.88,-1.95) -- (4.62,-1.95);
\draw[cA,line width=1.1pt] (2,-1.95) arc[start angle=180,end angle=0,x radius=1.25,y radius=0.75];
\draw[cA,line width=1.1pt] (2.5,-1.95) arc[start angle=180,end angle=0,x radius=0.25,y radius=0.15];
\draw[cA,line width=1.1pt] (3.5,-1.95) arc[start angle=180,end angle=0,x radius=0.25,y radius=0.15];
\fill (2,-1.95) circle (0.03);
\fill (2.5,-1.95) circle (0.03);
\fill (3,-1.95) circle (0.03);
\fill (3.5,-1.95) circle (0.03);
\fill (4,-1.95) circle (0.03);
\fill (4.5,-1.95) circle (0.03);
\node[font=\scriptsize,inner sep=0pt] at (2,-2.15) {1};
\node[font=\scriptsize,inner sep=0pt] at (2.5,-2.15) {2};
\node[font=\scriptsize,inner sep=0pt] at (3,-2.15) {3};
\node[font=\scriptsize,inner sep=0pt] at (3.5,-2.15) {4};
\node[font=\scriptsize,inner sep=0pt] at (4,-2.15) {5};
\node[font=\scriptsize,inner sep=0pt] at (4.5,-2.15) {6};
\draw[-{Stealth[length=5pt]},gray!60!black,line width=0.7pt] (4.95,-1.65) -- (6.38,-1.35) node[midway,above,sloped,font=\scriptsize] {捲回圓上};
\draw[gray!65,line width=0.5pt] (7.55,-1.2) circle (0.75);
\draw[cA,line width=1.1pt,line cap=round] (7.55,-0.45) -- (6.9,-0.825);
\draw[cA,line width=1.1pt,line cap=round] (8.2,-0.825) -- (8.2,-1.575);
\draw[cA,line width=1.1pt,line cap=round] (7.55,-1.95) -- (6.9,-1.575);
\fill[black] (7.55,-0.45) circle (0.035);
\fill[black] (8.2,-0.825) circle (0.035);
\fill[black] (8.2,-1.575) circle (0.035);
\fill[black] (7.55,-1.95) circle (0.035);
\fill[black] (6.9,-1.575) circle (0.035);
\fill[black] (6.9,-0.825) circle (0.035);
\node[font=\scriptsize] at (7.55,-0.21) {1};
\node[font=\scriptsize] at (8.407,-0.705) {2};
\node[font=\scriptsize] at (8.407,-1.695) {3};
\node[font=\scriptsize] at (7.55,-2.19) {4};
\node[font=\scriptsize] at (6.693,-1.695) {5};
\node[font=\scriptsize] at (6.693,-0.705) {6};
\end{tikzpicture}}

\newcommand{\FigMeightChain}{%
\begin{tikzpicture}
\fill[cTop!38] (-0.423,-0.651) -- (0.685,0.154) -- (0.423,-0.651) -- cycle;
\fill[cLeft!22] (-0.423,-0.651) -- (-0.685,0.154) -- (0,0.651) -- (0.685,0.154) -- cycle;
\draw[gray!70,line width=0.6pt,line join=round] (-0.423,-0.651) -- (-0.685,0.154) -- (0,0.651) -- (0.685,0.154) -- (0.423,-0.651) -- cycle;
\draw[cTop,line width=1.1pt,line cap=round] (-0.423,-0.651) -- (0.685,0.154);
\draw[cLeft,line width=1.1pt,line cap=round] (-0.685,0.154) -- (0.685,0.154);
\fill[black] (-0.423,-0.651) circle (0.03);
\fill[black] (-0.685,0.154) circle (0.03);
\fill[black] (0,0.651) circle (0.03);
\fill[black] (0.685,0.154) circle (0.03);
\fill[black] (0.423,-0.651) circle (0.03);
\node[font=\small] at (0,-1.05) {三角剖分};
\draw[cLeft,line width=0.8pt,line cap=round] (3.09,0.195) -- (2.67,-0.205);
\draw[cLeft,line width=0.8pt,line cap=round] (3.09,0.195) -- (3.3,-0.205);
\draw[cLeft,line width=0.8pt,line cap=round] (3.3,-0.205) -- (3.09,-0.605);
\draw[cLeft,line width=0.8pt,line cap=round] (3.3,-0.205) -- (3.51,-0.605);
\draw[cTop,line width=0.8pt,line cap=round] (3.3,0.595) -- (3.09,0.195);
\draw[cTop,line width=0.8pt,line cap=round] (3.3,0.595) -- (3.93,0.195);
\fill[cLeft] (3.09,0.195) circle (0.05);
\fill[cLeft] (3.3,-0.205) circle (0.05);
\fill[cTop] (3.3,0.595) circle (0.05);
\fill[cLeft] (2.63,-0.245) rectangle (2.71,-0.165);
\fill[cLeft] (3.05,-0.645) rectangle (3.13,-0.565);
\fill[cLeft] (3.47,-0.645) rectangle (3.55,-0.565);
\fill[cRight] (3.89,0.155) rectangle (3.97,0.235);
\node[font=\small] at (3.3,-1.05) {二元樹};
\draw[cTop,line width=1.2pt,line cap=round] (6.6,0.495) -- (6.6,-0.005);
\draw[cLeft,line width=1.2pt,line cap=round] (6.6,-0.005) -- (6.29,-0.505);
\draw[cLeft,line width=1.2pt,line cap=round] (6.6,-0.005) -- (6.91,-0.505);
\fill[white] (6.6,0.495) circle (0.065);
\draw[black,line width=0.6pt] (6.6,0.495) circle (0.065);
\fill[cTop] (6.6,-0.005) circle (0.055);
\fill[cLeft] (6.29,-0.505) circle (0.055);
\fill[cLeft] (6.91,-0.505) circle (0.055);
\node[font=\small] at (6.6,-1.05) {平面樹};
\draw[gray!55,line width=0.4pt] (9.12,-0.26) -- (10.68,-0.26);
\draw[cTop,line width=1.3pt,line cap=round] (9.12,-0.26) -- (9.38,0);
\draw[cLeft,line width=1.3pt,line cap=round] (9.38,0) -- (9.64,0.26);
\draw[cLeft,line width=1.3pt,line cap=round] (9.64,0.26) -- (9.9,0);
\draw[cLeft,line width=1.3pt,line cap=round] (9.9,0) -- (10.16,0.26);
\draw[cLeft,line width=1.3pt,line cap=round] (10.16,0.26) -- (10.42,0);
\draw[cTop,line width=1.3pt,line cap=round] (10.42,0) -- (10.68,-0.26);
\node[font=\small] at (9.9,-1.05) {Dyck 路徑};
\draw[gray!65,line width=0.5pt] (13.2,-0.015) circle (0.62);
\draw[cTop,line width=1.1pt,line cap=round] (13.2,0.605) -- (12.663,0.295);
\draw[cLeft,line width=1.1pt,line cap=round] (13.737,0.295) -- (13.737,-0.325);
\draw[cLeft,line width=1.1pt,line cap=round] (13.2,-0.635) -- (12.663,-0.325);
\fill[black] (13.2,0.605) circle (0.035);
\fill[black] (13.737,0.295) circle (0.035);
\fill[black] (13.737,-0.325) circle (0.035);
\fill[black] (13.2,-0.635) circle (0.035);
\fill[black] (12.663,-0.325) circle (0.035);
\fill[black] (12.663,0.295) circle (0.035);
\fill[cTop] (13.2,0.605) circle (0.065);
\node[font=\small] at (13.2,-1.05) {不交叉配對};
\draw[{Stealth[length=4pt]}-{Stealth[length=4pt]},gray!60!black,line width=0.7pt] (0.95,0) -- (2.35,0) node[midway,above,font=\scriptsize,text=gray!40!black] {對偶圖};
\draw[{Stealth[length=4pt]}-{Stealth[length=4pt]},gray!60!black,line width=0.7pt] (4.25,0) -- (5.65,0) node[midway,above,font=\scriptsize,text=gray!40!black] {小孩／兄弟};
\draw[{Stealth[length=4pt]}-{Stealth[length=4pt]},gray!60!black,line width=0.7pt] (7.55,0) -- (8.95,0) node[midway,above,font=\scriptsize,text=gray!40!black] {輪廓走法};
\draw[{Stealth[length=4pt]}-{Stealth[length=4pt]},gray!60!black,line width=0.7pt] (10.85,0) -- (12.25,0) node[midway,above,font=\scriptsize,text=gray!40!black] {U 配 D};
\end{tikzpicture}}

% =====================================================================
% 附錄 A 第 2 段：A.5–A.8（動畫 M05–M08）
% 圖在 mit/figs_part2.tex，由 mit/make_figs_part2.py 產生（含數字檢查）。
% =====================================================================

% ---------------------------------------------------------------------
\subsection{平面樹}\label{sec:m05}

第 1.5 節的滿二元樹，每個內部節點恰有左、右兩個小孩。\textbf{平面樹}（plane tree，也叫有序樹）放寬這個限制
\mitref{L3}：它有一個\textbf{根}，每個頂點可以有任意多個小孩，而且小孩\textbf{由左到右排好順序}。
$n+1$ 個頂點的樹恰有 $n$ 條邊，所以用邊數 $n$ 當作大小。
順序是平面樹的一部分：圖~\ref{fig:m05-trees} 左邊的兩棵只差在根的兩個小孩左右交換，
不管順序時是同一棵樹，作為平面樹卻是不同的兩棵。

\begin{figure}[H]
\centering
\FigMfiveTrees
\caption{左：交換根的兩個小孩（連同下面的部分），得到不同的平面樹。
右：$n=0,1,2,3$ 條邊的所有平面樹；虛線框是左邊那兩棵（第 2、3 棵），\textcolor{cTop}{橘框}是貫穿的例子。}
\label{fig:m05-trees}
\end{figure}

\begin{thmbox}[平面樹]
$n$ 條邊（$n+1$ 個頂點）的平面樹恰有 $C_n$ 棵：$n=0,1,2,3,4$ 時是 $1,1,2,5,14$。
\end{thmbox}

\textbf{根＋左＋右。}設 $n\ge1$。看根的\textbf{第一個小孩} $v$，把根到 $v$ 的這條邊當作「根＋左＋右」的根：
$v$ 和它下面的部分是一棵平面樹，設有 $k$ 條邊（左）；原來的根連同其餘的小孩（和它們下面的部分）也是一棵平面樹，
有 $n-1-k$ 條邊（右）。反過來，把左邊那棵掛到右邊那棵的根底下、排在最左邊，就拼回原來的樹。兩部分互不干擾，
所以 $n$ 條邊的平面樹個數 $P_n$ 滿足
\[
P_n=\sum_{k=0}^{n-1}P_k\,P_{n-1-k}\quad(n\ge1),\qquad P_0=1,
\]
正是遞迴式 \eqref{eq:segner}，因此 $P_n=C_n$。
圖~\ref{fig:m05-trees} 中 $n=3$ 那一列依序對應圖~\ref{fig:onenumber} 的五欄（換法見第~\ref{sec:m07} 節）。

% ---------------------------------------------------------------------
\subsection{弧圖：不交叉與不巢狀}\label{sec:m06}

把第 1.10 節的圓在點 $2n$ 與點 $1$ 之間剪開、攤平成一條直線，點依序變成直線上的 $1,2,\ldots,2n$，
弦變成直線上方的\textbf{弧}，這叫配對的\textbf{弧圖} \mitref{L3}。兩條弧的端點記為 $i<j<k<l$，配法只有三種：
$(i,j),(k,l)$ \textbf{分開}；$(i,k),(j,l)$ \textcolor{cBad}{\textbf{交叉}}；
$(i,l),(j,k)$ \textcolor{violet!75!black}{\textbf{巢狀}}，也就是一條包住另一條（圖~\ref{fig:m06-cut}）。
圓上兩條弦相交，正好是端點在圓周上交錯，剪開後就是交叉的情形，所以圓上的不交叉配對就是沒有交叉的弧圖。
貫穿的例子 $(1,6),(2,3),(4,5)$ 沒有交叉，但 $(1,6)$ 包住了另外兩條：有巢狀。

\begin{figure}[H]
\centering
\FigMsixCut
\caption{左：在點 6 與點 1 之間剪開，圓上的配對 $(1,6),(2,3),(4,5)$ 攤平成直線上的弧圖。
右：兩條弧的三種關係。}
\label{fig:m06-cut}
\end{figure}

\begin{thmbox}[不交叉與不巢狀]
直線上 $2n$ 個點的完美配對中，\textbf{不交叉}的有 $C_n$ 種，\textbf{不巢狀}的也有 $C_n$ 種。
\end{thmbox}

$n=3$：6 個點的完美配對共 $5\times3\times1=15$ 種（點 1 有 5 個對象可選，剩下最左邊的點有 3 個，最後兩點只有 1 種配法）。
其中不交叉的 5 種、不巢狀的 5 種，$(1,2),(3,4),(5,6)$ 兩者都是，其餘 6 種既有交叉也有巢狀（圖~\ref{fig:m06-count} 左）。

\textbf{證明：兩種都對應到 Dyck 路徑。}把每條弧的左端記 U、右端記 D。每個 D 所在那條弧的 U 都在它前面，
所以讀到任何位置 D 都不比 U 多：得到一條 Dyck 路徑。反過來，給一條 Dyck 路徑，要把每個 D 配給它前面一個還沒配的 U。
設位置 $k$ 的 D 配給位置 $i$ 的 U：若 $i$ 與 $k$ 之間還有沒配的 U（位置 $j$），它以後配到某個 $l>k$，就有交叉 $(i,k),(j,l)$；
若 $i$ 之前還有沒配的 U（位置 $h$），它以後配到某個 $l>k$，就有巢狀：$(h,l)$ 包住 $(i,k)$。
反之，每個交叉、每個巢狀都是這樣產生的（看兩條弧裡較早出現的右端）。所以

\begin{keybox}[同一條 Dyck 路徑，兩種配法]
\textbf{不交叉} $\iff$ 每個 D 都配給\textbf{最近}的、還沒配的 U（後進先出，像堆疊）；\\
\textbf{不巢狀} $\iff$ 每個 D 都配給\textbf{最早}的、還沒配的 U（先進先出，像佇列）。\\
兩條規則對每一條 Dyck 路徑都各給出唯一的配對，所以兩種配對都和 Dyck 路徑一一對應，各有 $C_n$ 種。
（堆疊與佇列在第~\ref{sec:m09} 節還會再出現。）
\end{keybox}

\begin{figure}[H]
\centering
\FigMsixCount
\caption{左：6 個點的 15 種完美配對，每一欄的點 1 配同一個點；\textcolor{cA}{藍框}是不交叉（5 種），
\textcolor{green!50!black}{綠色虛線框}是不巢狀（5 種）。右：同一個 Dyck 字（左端 U、右端 D）
對應一個不交叉配對和一個不巢狀配對；\textcolor{cTop}{橘框}是 UUDUDD：$(1,6),(2,3),(4,5)$ 與 $(1,3),(2,5),(4,6)$。}
\label{fig:m06-count}
\end{figure}

% ---------------------------------------------------------------------
\subsection{三角剖分的對偶樹}\label{sec:m07}

第 1.11 節經由「根＋左＋右」的遞迴得到三角剖分與滿二元樹的一一對應；這棵樹其實就是三角剖分的\textbf{對偶圖}
\mitref{L3}。凸 $(n+2)$ 邊形的頂點依序編號為 $1,\ldots,n+2$，底邊是 $\{1,n+2\}$：
\begin{enumerate}[leftmargin=1.6em]
  \item 每個三角形放一個節點；含底邊的三角形是\textcolor{cTop}{\textbf{根}}。
  \item 兩個三角形共用一條對角線 $\to$ 兩個節點連一條邊。
  \item 其他每條多邊形的邊（底邊除外）外面放一片\textbf{葉子}，連到含這條邊的三角形。
\end{enumerate}
三角形 $\{a,k,b\}$（$a<k<b$）中，$\{a,b\}$ 是朝向底邊的那條邊（根三角形的 $\{a,b\}$ 就是底邊）；
越過 $\{a,k\}$ 的是\textbf{左小孩}，越過 $\{k,b\}$ 的是\textbf{右小孩}。
$n$ 個三角形、$n-1$ 條對角線、$n+1$ 條非底邊，所以得到 $n$ 個內部節點、$n+1$ 片葉子的滿二元樹。

\textbf{貫穿的例子}（圖~\ref{fig:m07-dual} 左）。五邊形、對角線 $\{1,4\},\{2,4\}$，三個三角形是
\textcolor{cTop}{$\mathrm{A}=\{1,4,5\}$}（根）、$\mathrm{B}=\{1,2,4\}$、$\mathrm{C}=\{2,3,4\}$。
A 越過 $\{1,4\}$ 是 B（左小孩），越過 $\{4,5\}$ 是葉子；B 越過 $\{1,2\}$ 是葉子，越過 $\{2,4\}$ 是 C（右小孩）；
C 的 $\{2,3\},\{3,4\}$ 都是葉子。得到滿二元樹 $((\bullet,(\bullet,\bullet)),\bullet)$，4 片葉子對應 4 條非底邊。

\begin{figure}[H]
\centering
\FigMsevenDual
\caption{左：五邊形的三角剖分與它的對偶圖（灰色的連線，灰色方塊是葉子），\textcolor{cTop}{橘色}三角形含底邊，是根；
拿掉多邊形、整理成標準畫法，顏色同圖~\ref{fig:onenumber}（\textcolor{cTop}{根}、\textcolor{cLeft}{左子樹}、
\textcolor{cRight}{右子樹}）。右：去掉葉子，\textcolor{cLeft}{左小孩}變成第一個小孩、\textcolor{cRight}{右小孩}變成下一個兄弟，
再加一個新的根。}
\label{fig:m07-dual}
\end{figure}

\textbf{為什麼是一一對應。}根三角形 $\{1,k,n+2\}$ 把多邊形切成左邊（頂點 $1,\ldots,k$）與右邊（頂點 $k,\ldots,n+2$），
兩塊的對偶樹正是根的左、右子樹——這就是第 1.3 節的「固定底邊」。反過來，若左子樹有 $j$ 個內部節點，
根三角形的第三個頂點就是 $k=j+2$，再往兩塊遞迴，就切回唯一的三角剖分。
例子中左子樹有 B、C 兩個節點，$k=4$：左邊是\textcolor{cLeft}{四邊形 $\{1,2,3,4\}$}，
右邊只剩\textcolor{cRight}{邊 $\{4,5\}$}（一片葉子）。

\textbf{二元樹 $\to$ 平面樹。}葉子只是空位，拿掉後剩下 $n$ 個節點的二元樹（第 1.5 節的另一種說法）。規則是
\[
\textcolor{cLeft}{\text{左小孩}}\ \longmapsto\ \text{第一個小孩},\qquad
\textcolor{cRight}{\text{右小孩}}\ \longmapsto\ \text{下一個兄弟（同一個父親，排在右邊）},
\]
再加一個新的根，二元樹的根是它的第一個小孩，得到 $n+1$ 個頂點的平面樹。
反過來，把每個非根頂點的第一個小孩當左小孩、下一個兄弟當右小孩，再拿掉根，就回到二元樹。
例子（圖~\ref{fig:m07-dual} 右）：B 是 A 的左小孩 $\to$ A 的第一個小孩；C 是 B 的右小孩 $\to$ B 的下一個兄弟，
也就是 A 的第二個小孩；A 沒有右小孩，所以新的根只有 A 一個小孩。
得到「根 $\to$ A $\to$ B、C」，正是圖~\ref{fig:m05-trees} 的第 4 棵。
這條規則就是第~\ref{sec:m05} 節的拆法：第一個小孩 $v$ 是根，$v$ 下面的部分是左子樹，其餘的部分是右子樹。

% ---------------------------------------------------------------------
\subsection{DFS：樹 → Dyck → 配對}\label{sec:m08}

從根出發，沿著平面樹的輪廓走一圈：先走最左邊的小孩，它的子樹走完回到父親，再走下一個小孩
（\textbf{深度優先搜尋}，DFS）\mitref{L3}。往下走一條邊記 U，往上走一條邊記 D。
每條邊恰好走兩次（下去一次、上來一次），所以 $n$ 條邊給出 $n$ 個 U、$n$ 個 D；
走到任何時候，高度（U 數減 D 數）就是所在頂點的深度，不會小於 0，最後回到根是 0。
所以得到一條半長 $n$ 的 Dyck 路徑。貫穿的例子（圖~\ref{fig:m08-walk}）：
\[
\text{根}\to\mathrm{A}\to\mathrm{B}\to\mathrm{A}\to\mathrm{C}\to\mathrm{A}\to\text{根}
\quad\Longrightarrow\quad \mathrm{UUDUDD},\qquad \text{高度}\ 0,1,2,1,2,1,0 .
\]
反過來，讀一條 Dyck 路徑：U 表示在目前的頂點下面長出一個新的小孩（排在已有的小孩右邊）並走下去，D 表示走回父親，
就重建出原來的平面樹。所以 DFS 是一一對應，也是第~\ref{sec:m05} 節定理的另一個證明。
它和第 1.7 節的「第一次回到地面」一致：第一步走進根的第一個小孩 $v$，第一次回到地面是從 $v$ 走回根，
中間是 $v$ 下面的部分（左），之後是其餘的部分（右）。

\begin{figure}[H]
\centering
\FigMeightWalk
\caption{左：沿平面樹的輪廓走一圈，邊旁的數字是第幾步（$\downarrow$ 往下、$\uparrow$ 往上）。
中：記下的 Dyck 路徑 UUDUDD，以及同一條邊的 U、D 連成的弧。右：捲回圓上，得到不交叉配對 $(1,6),(2,3),(4,5)$。}
\label{fig:m08-walk}
\end{figure}

\textbf{U 配 D。}同一條邊「下去」的 U 和「上來」的 D 配成一對。例子中根--A 是第 1、6 步，A--B 是第 2、3 步，
A--C 是第 4、5 步，得到配對 $(1,6),(2,3),(4,5)$。一條邊的 U 與 D 之間，走的正好是這條邊下面的子樹；
兩條邊若一條在另一條下面，弧就是巢狀，否則兩段走訪不重疊，弧就分開。所以弧\textbf{不會交叉}，
每個 D 配的都是最近的、還沒配的 U（第~\ref{sec:m06} 節的後進先出）。
把直線的兩端接起來、捲回圓上（第~\ref{sec:m06} 節的反向），弧變成弦，就是第 1.10 節圓上的不交叉配對。

\begin{figure}[H]
\centering
\FigMeightChain
\caption{貫穿的例子的五種樣子，箭頭上是對應的規則。顏色依第 1 節的根／左／右：根三角形、樹根、平面樹的 A、
第一個 U 與最後一個 D、弦 $(1,6)$ 是\textcolor{cTop}{根}；B、C 那一部分是\textcolor{cLeft}{左}；右是空的。}
\label{fig:m08-chain}
\end{figure}

\begin{keybox}[一條鏈]
三角剖分 $\leftrightarrow$ 二元樹 $\leftrightarrow$ 平面樹 $\leftrightarrow$ Dyck 路徑 $\leftrightarrow$ 不交叉配對：
每一步都是直接構造出來的一一對應，所以這五種物件的個數都是 $C_n$（$n=3$ 時都是 5）。
合起來的對應和第 1.11 節經由「根＋左＋右」得到的完全一致（圖~\ref{fig:onenumber} 的同一欄）；
程式對 $n\le6$ 的每一棵樹都檢查過。
\end{keybox}
